\subsection{Course Exercises}

\begin{framed}
\textbf{Golden rule}\\
Always \textbf{attempt the exercise yourself} before looking at the solution. Struggling with a problem is where learning happens!
\end{framed}


\bigskip
\centerline{\rule{13cm}{0.4pt}}
\bigskip

\paragraph{Structure}

Each week's exercise folder contains:

\begin{itemize}
\item \textbf{Main exercise notebook} --- problems with guidance and hints
\item \textbf{\texttt{solutions/} sub-folder} --- reference implementations to review after your attempt
\end{itemize}

\begin{verbatim}
exercises/
├── week1_python_fundamentals/
│   ├── exercise.ipynb
│   └── solutions/
│       └── solution.ipynb
├── week2_resipy/
│   ├── exercise.ipynb
│   └── solutions/
│       └── solution.ipynb
├── week3_emagpy/
│   ├── exercise.ipynb
│   └── solutions/
│       └── solution.ipynb
└── week4_case_study/
    ├── exercise.ipynb
    └── solutions/
        └── solution.ipynb
\end{verbatim}


\bigskip
\centerline{\rule{13cm}{0.4pt}}
\bigskip

\paragraph{Weekly Topics}


\bigskip
\centerline{\rule{13cm}{0.4pt}}
\bigskip

\paragraph{Guidelines}

\begin{table}
\centering
\begin{tabular}{p{\dimexpr 0.500\linewidth-2\tabcolsep}p{\dimexpr 0.500\linewidth-2\tabcolsep}}
\toprule
\subsubsection{}

 & Guideline \\
\hline
1 & Attempt exercises \textbf{before} viewing solutions \\
2 & Experiment --- change parameters and observe what happens \\
3 & Compare your approach with the provided solution; both can be correct! \\
4 & Ask questions in the course discussion if you are stuck for more than 20 minutes \\
\bottomrule
\end{tabular}
\end{table}


\bigskip
\centerline{\rule{13cm}{0.4pt}}
\bigskip

\begin{framed}
\textbf{Tip}\\
Use the \textbf{Jupyter Lab} split-screen view (\texttt{Ctrl + Shift + D}) to keep the exercise and solution side by side when reviewing.
\end{framed}